
\documentclass[12pt,a4paper]{article}

% Packages
\usepackage[margin=1in]{geometry}
\usepackage{amsmath,amssymb}
\usepackage{graphicx}
\usepackage{booktabs}
\usepackage{hyperref}
\usepackage{setspace}
\usepackage{xcolor}

\onehalfspacing

% Document information
\title{\textbf{Technical Writing and Scientific Typesetting}\\
\large A Study of Digital Documentation Tools}
\author{Sudeep Abraham Ajith\\
Department of Computer Science and Engineering\\
Malaviya National Institute of Technology Jaipur}
\date{November 2026}

\begin{document}

% Title page
\maketitle
\thispagestyle{empty}
\newpage

% Abstract
\begin{abstract}
Technical writing plays an important role in communicating
complex information clearly and effectively. Scientific
typesetting tools such as LaTeX provide structured methods
for preparing professional reports, academic documents,
mathematical expressions, tables, and references.

This report examines the fundamental features of LaTeX,
its document structure, mathematical typesetting
capabilities, and its applications in technical
documentation.
\end{abstract}

\newpage
\tableofcontents
\newpage

\section{Introduction}

Technical writing is the process of presenting
specialised information in a clear, organised,
and accurate manner. It is widely used in engineering,
science, research, and software development.

LaTeX is a document preparation system particularly
suited to scientific and technical publications.
Unlike conventional word processors, it separates
document content from formatting instructions.

\section{Objectives}

The main objectives of this report are:

\begin{itemize}
    \item To understand the purpose of LaTeX.
    \item To examine the basic structure of a LaTeX document.
    \item To demonstrate mathematical typesetting.
    \item To understand the creation of tables and lists.
    \item To explore the role of LaTeX in technical reports.
\end{itemize}

\section{Structure of a LaTeX Document}

A typical LaTeX document consists of two major parts:

\begin{enumerate}
    \item \textbf{Preamble:} Contains the document class,
    packages, and formatting settings.
    \item \textbf{Document body:} Contains the actual
    text, headings, equations, tables, and references.
\end{enumerate}

The following example illustrates the basic structure:

\begin{verbatim}
\documentclass{article}
\begin{document}
Hello, World!
\end{document}
\end{verbatim}

\section{Mathematical Typesetting}

One of the major advantages of LaTeX is its ability
to represent complex mathematical expressions.

For example, the quadratic formula is:

\begin{equation}
x = \frac{-b \pm \sqrt{b^2 - 4ac}}{2a}
\end{equation}

The summation notation can be represented as:

\begin{equation}
S = \sum_{i=1}^{n} x_i
\end{equation}

The probability density function of a normal
distribution is:

\begin{equation}
f(x) =
\frac{1}{\sigma\sqrt{2\pi}}
e^{-\frac{(x-\mu)^2}{2\sigma^2}}
\end{equation}

\section{Tables and Document Organisation}

LaTeX allows information to be presented in
structured tables.

\begin{table}[h]
\centering
\caption{Common LaTeX Commands}
\begin{tabular}{ll}
\toprule
\textbf{Command} & \textbf{Purpose} \\
\midrule
\texttt{\textbackslash section} & Creates a section \\
\texttt{\textbackslash subsection} & Creates a subsection \\
\texttt{\textbackslash textbf} & Applies bold formatting \\
\texttt{\textbackslash includegraphics} & Inserts an image \\
\texttt{\textbackslash tableofcontents} & Creates a contents page \\
\bottomrule
\end{tabular}
\end{table}

\section{Advantages of LaTeX}

The major advantages of LaTeX include:

\begin{itemize}
    \item Consistent and professional document formatting.
    \item High-quality mathematical typesetting.
    \item Automatic numbering of sections and equations.
    \item Support for tables, figures, and references.
    \item Suitability for large technical documents.
\end{itemize}

\section{Conclusion}

LaTeX is a powerful tool for preparing structured
technical and scientific documents. Its mathematical
typesetting capabilities, consistent formatting,
and flexible document organisation make it useful
for academic reports and research publications.

Understanding LaTeX provides a foundation for
creating professional technical documentation.

\begin{thebibliography}{9}

\bibitem{lamport}
Leslie Lamport,
\textit{LaTeX: A Document Preparation System},
2nd Edition, Addison-Wesley, 1994.

\bibitem{overleaf}
Overleaf,
\textit{Learn LaTeX Documentation}.
\url{https://www.overleaf.com/learn}

\end{thebibliography}

\end{document}
